\documentclass[sigconf]{acmart}

% Remove the ACM copyright/conference boilerplate for a course report.
\settopmatter{printacmref=false}
\setcopyright{none}
\renewcommand\footnotetextcopyrightpermission[1]{}
\acmConference[PRI 2026/27]{Information Processing and Retrieval}{FEUP}{Porto, Portugal}

\usepackage{xcolor}
% Visible reminders for things the team still has to confirm. Delete before submission.
\newcommand{\todo}[1]{\textcolor{red}{\textbf{[TODO: #1]}}}
% Include a figure if it exists, otherwise draw a placeholder so the report still compiles.
\newcommand{\figorbox}[2]{\IfFileExists{#1}{\includegraphics[width=#2]{#1}}{\fbox{\parbox{#2}{\centering\vspace{1.5em}missing: \texttt{#1}\vspace{1.5em}}}}}

\begin{document}

% ==================================================
% TITLE
% ==================================================

\title{Pok\'emon Cards: Collection, Processing and Characterization}

% ==================================================
% AUTHORS
% ==================================================

\author{Afonso Saraiva}
\affiliation{%
  \institution{Faculty of Engineering, University of Porto}
  \city{Porto}
  \country{Portugal}
}
\email{up202304461@up.pt}

\author{Inês Francisco}
\affiliation{%
  \institution{Faculty of Engineering, University of Porto}
  \city{Porto}
  \country{Portugal}
}
\email{up202304726@up.pt}

\author{Leonardo Carvalho}
\affiliation{%
  \institution{Faculty of Engineering, University of Porto}
  \city{Porto}
  \country{Portugal}
}
\email{up202307152@up.pt}

\author{Miguel Pereira}
\affiliation{%
  \institution{Faculty of Engineering, University of Porto}
  \city{Porto}
  \country{Portugal}
}
\email{up202304387@up.pt}

\renewcommand{\shortauthors}{Saraiva et al.}

% ==================================================
% ABSTRACT
% ==================================================

\begin{abstract}
The Pok\'emon Trading Card Game (TCG) has printed tens of thousands of cards
since 1996, and players and collectors routinely need to find cards by what
they \emph{do}, who drew them, or where they came from, rather than by exact
name. This report describes the first milestone of a search system for that
domain. We collected card, species, illustrator and expansion pages from
Bulbapedia and joined them with a TCGplayer market-price snapshot obtained
through TCGCSV. A Python pipeline parses the wiki pages, groups reprints into
13{,}972 unique cards with 18{,}420 prints, links each card to 1{,}027
Pok\'emon species, 427 illustrators and 221 expansions, and attaches prices to
94\% of the cards. The resulting documents average 89 words of free text
(1.25M tokens in total), and carry 20{,}327 structured attacks with
Energy cost and damage. We characterize the collection, including field
completeness, release history, the growth of card HP over time, price
distributions and the vocabulary of rules text, and define five information
needs that will drive the retrieval milestones.
\end{abstract}

\keywords{Pok\'emon TCG, data collection, data preparation, web scraping,
information retrieval}

\maketitle

% ==================================================
% 1. INTRODUCTION
% ==================================================

\section{Introduction}

A Pok\'emon TCG card combines short structured attributes (type, HP,
weakness, retreat cost) with natural-language rules text describing
its attacks and abilities. Around each card there is a body of
encyclopaedic text: how and where it was released, trivia about its
artwork, the biology of the depicted species and the career of its
illustrator. Official card databases offer filters on the structured
fields but do not rank cards by how well their text answers a need such
as ``cards that recover Energy from the discard pile''.

The goal of this project is a search system that returns a ranked list of
cards for such needs. This first milestone covers selecting and assessing
data sources, building a reproducible collection pipeline, defining the
conceptual model and the document unit, characterizing the collection,
and formulating the information needs that later milestones will evaluate.

% ==================================================
% 2. DATA SOURCES
% ==================================================

\section{Data Sources}

\paragraph{Bulbapedia~\cite{bulbapedia}.}
A community-maintained Pok\'emon encyclopaedia built on MediaWiki, with
a dedicated page for almost every TCG card, species, illustrator and
expansion. Pages contain structured infoboxes (wikitext templates) and
long prose sections. Bulbapedia is the most complete public reference for
the TCG and is curated by an editorial team, which we considered sufficient
authority for card text and release history. Content is licensed under
CC BY-NC-SA 2.5~\cite{ccbyncsa}, which permits our non-commercial academic
use with attribution. Pages were retrieved as raw wikitext through the
MediaWiki API~\cite{mediawikiapi} rather than by scraping rendered HTML,
with an identifying User-Agent, a 2-second delay between requests and the
API's \texttt{maxlag} back-off. \todo{crawl date}

\paragraph{TCGCSV~\cite{tcgcsv}.}
A public, daily mirror of the TCGplayer~\cite{tcgplayer} catalogue and price
guide, organised as \emph{groups} (product lines/sets) and \emph{products},
each with market, low, mid and direct-low prices per finish (normal,
holofoil, reverse holofoil). TCGplayer is the largest card marketplace in
North America, so its prices are a reasonable market signal. We used the
snapshot of 2026-10-11 (213 groups, 28{,}054 products). Prices remain
TCGplayer's data and are used here for non-commercial research only.

\paragraph{Sources evaluated and discarded.}
We also built extractors for the Pok\'emon TCG API~\cite{pokemontcgapi},
Trainer Hill~\cite{trainerhill} and Pok\'emonMeta~\cite{pokemonmeta}.
The TCG API's free tier (100 requests/day) did not allow a complete
capture. Trainer Hill exposes only archetype-level win-rate matchups
(778 pairs across 45 archetypes, 2020--2026) without decklists or card
identifiers, so it cannot be linked to individual cards. Pok\'emonMeta
covers \emph{TCG Pocket}, a different game with its own card pool.
Their extractors are kept on a separate branch of the repository but
are not part of the collection.

\begin{table}[t]
  \caption{Final collection (Bulbapedia + TCGCSV snapshot).}
  \label{tab:collection}
  \small
  \begin{tabular}{lrl}
    \toprule
    Entity & Records & Main text fields (median words) \\
    \midrule
    Card & 13{,}972 & rules (24), release info (52), trivia (24) \\
    Print & 18{,}420 & -- \\
    Price row & 33{,}835 & -- \\
    Pok\'emon species & 1{,}027 & biology (267), 10{,}700 Pok\'edex entries \\
    Artist & 427 & description (50) \\
    Expansion & 221 & information (194), blurb \\
    \bottomrule
  \end{tabular}
\end{table}

% ==================================================
% 3. COLLECTION AND PREPARATION
% ==================================================

\section{Collection and Preparation}

\subsection{Pipeline}

Figure~\ref{fig:pipeline} shows the pipeline, implemented in Python in
four stages (\texttt{fetching/}, \texttt{extracting/}, \texttt{processing/}
and a quality check).

\paragraph{Fetching.}
Starting from Bulbapedia's list of TCG expansions, the crawler reads each
expansion page, keeps the links that follow the card-page naming pattern
(\textit{Name (Expansion Number)}) and downloads their wikitext in batches
of up to 50 titles per request, resolving redirects. Species and illustrator
pages are fetched in a second pass, using the names referenced by the
extracted cards. For prices, the TCGCSV group list and one
\texttt{ProductsAndPrices.csv} per group are downloaded. All raw responses
are cached on disk, so interrupted runs resume without re-downloading.

\paragraph{Extracting.}
Wikitext is parsed with \texttt{mwparserfromhell}. For each card we keep
the infobox fields (name, species, stage, evolution, type, HP, weakness,
resistance, retreat cost, trainer class), the structured card text (each
attack or ability with its Energy cost, damage and effect), Pok\'edex
entries, the image gallery, the \emph{Release information} and
\emph{Trivia} sections, and the page categories. Japanese-only fields and
prints are dropped.

\paragraph{Processing.}
\texttt{process\_cards.py} assigns stable slug identifiers
(e.g.\ \texttt{card:abra-151-63}), joins each print to its expansion,
illustrator and image, and links cards to species, illustrators and
expansions by normalized name (case, accents and spacing ignored). Names
without a page become stub records (13 artists, 1 species, several trainer
kits) so that no link is dangling. Prints missing an expansion, illustrator
or image are dropped, as are cards without an English print. The card text
is also flattened into a single \texttt{rules\_text} field for full-text
search, and the remaining categories become tags.
\texttt{merge\_prices.py} matches each TCGplayer product to a print by
cleaned card name and number; a TCGCSV group is mapped to the expansion
most of its products match (at least 3 votes; 181 of 213 groups), which
breaks ties and enables a number-only fallback. 21{,}696 products matched,
6{,}290 did not (mostly jumbo cards, world-championship decks and promos)
and 68 were ambiguous and left out rather than guessed. Each step writes a
report (\texttt{links\_report.json}, \texttt{price\_report.json}) listing
collisions, stubs and unmatched records, which we used for quality
assessment.

\todo{describe the final cleaning step and the single merged document file:
normalize the 46 raw \texttt{kind} values (e.g.\ \texttt{svex},
\texttt{mega ex}, \texttt{megaex}) into card category + mechanic, strip
whitespace in artist names, drop invalid stub artists such as
``(This card cannot be used at official tournaments.)''}

\paragraph{Reproducibility.}
The whole pipeline runs with \texttt{make pipeline} (\texttt{uv run main.py})\nand the characterization with \texttt{make eda}; dependencies are
pinned in \texttt{pyproject.toml}/\texttt{uv.lock}, and
\texttt{eda/characterize.py} regenerates every statistic and figure in
Section~\ref{sec:char}. \todo{fix main.py imports (extract\_prices.py
missing, build\_quality\_report undefined) before submission}

\begin{figure*}[t]
  \centering
  \figorbox{figures/pipeline.pdf}{\textwidth}
  \caption{Data collection and preparation pipeline.}
  \label{fig:pipeline}
\end{figure*}

\subsection{Conceptual Data Model}

Figure~\ref{fig:uml} shows the conceptual model. \emph{Card} is the central
entity and the retrieval document. A card has one or more \emph{Prints};
each print belongs to one \emph{Expansion}, is illustrated by one
\emph{Artist} and may have several \emph{PriceSnapshots} (one per product
finish). A card has zero or more \emph{CardText} entries: attacks (with
Energy cost, damage and effect), abilities and trainer effects. A card
depicts zero or more \emph{Pok\'emon species} (zero for
Trainer and Energy cards, several for tag-team cards), and each species
has \emph{Pok\'edex entries} attached to the games in which they appear.

\begin{figure}[t]
  \centering
  \figorbox{figures/uml.pdf}{\columnwidth}
  \caption{Conceptual data model.}
  \label{fig:uml}
\end{figure}

\subsection{Document Definition}

Each document in the final collection is one card, identified by its slug
(e.g.\ \texttt{card:abra-151-63}), with a title (name plus first print),
structured fields for filtering and its own text fields (rules text, release
information, trivia), enriched with denormalized text from the linked
entities: species biology and Pok\'edex entries, illustrator description and
expansion information. This lets a query about an illustrator or a species'
habitat match the card even when the card text does not mention it.

% ==================================================
% 4. DATA CHARACTERIZATION
% ==================================================

\section{Data Characterization}
\label{sec:char}

\paragraph{Composition.}
Of the 13{,}972 cards, 12{,}542 (89.8\%) are Pok\'emon cards, 1{,}329 are
Trainer cards (Item, Supporter, Stadium, Tool) and 101 are other cards,
mostly Special Energy (87); basic Energy cards are not in the collection. Types are balanced
(Water 1{,}816, Grass 1{,}799, Psychic 1{,}691, Colorless 1{,}646) with Dragon
(317) and Fairy (172) smaller, as these types were introduced later. 80\%
of cards have a single print; the most reprinted card has 17.

\paragraph{Completeness and quality.}
Identifiers, names, rules text, artists and expansions are complete.
Pok\'emon-only attributes are missing for non-Pok\'emon cards as expected
(HP 10.2\%, type 9.5\% missing); resistance is absent on 74.5\% of cards
because many Pok\'emon have none. Among text fields, trivia is present on
92.6\% of cards but release information only on 43.1\%. 94\% of
cards (13{,}120) have a market price; the median minimum price is US\$0.36,
the 95th percentile US\$34 and the maximum US\$3{,}000 (Charizard,
Skyridge~146).

\paragraph{Attacks and abilities.}
The 13{,}972 cards carry 26{,}675 card-text entries: 20{,}327 attacks
and 2{,}642 abilities (including the older Pok\'e-Power and Pok\'e-Body).
98.8\% of attacks have an Energy cost, most often of one to three Energy
(7{,}953, 6{,}093 and 4{,}886 attacks), and 79\% have a printed damage value
(median 30, maximum 330). In 3{,}826 of these the damage is variable,
written as ``120+'', ``20$\times$'' or ``50$-$'' and explained in the effect
text.

\paragraph{Variables over time.}
Release years were parsed from expansion pages (125 of 221 expansions;
12{,}688 cards dated). Figure~\ref{fig:years} shows the number of cards per
year, which grows sharply after 2016. Figure~\ref{fig:hp} shows HP over time:
the median HP of Pok\'emon cards rose from 60--70 in 1999--2002 to 110 in
2026, and the maximum from 120 to 380, which reflects the ``power creep''
widely discussed by players. Attack damage shows the same trend
(Figure~\ref{fig:dmg}): the 90th percentile rose from 40 in 2000 to 160 in
2026.
\todo{expansions after 2024 are under-dated by the year parser, so the
2025 bar is incomplete; fix before the final version}

\begin{figure}[t]
  \centering
  \figorbox{figures/cards_per_year.pdf}{\columnwidth}
  \caption{Cards per release year (year of earliest print).}
  \label{fig:years}
\end{figure}

\begin{figure}[t]
  \centering
  \figorbox{figures/hp_over_time.pdf}{\columnwidth}
  \caption{Median and maximum HP of Pok\'emon cards per release year.}
  \label{fig:hp}
\end{figure}

\paragraph{Text analysis.}
A card document has a median of 69 words (mean 89, max 2{,}839), giving
1.25M tokens and a vocabulary of 12{,}147 terms, 3{,}785 of which occur once.
The most frequent rules-text terms after stop-word removal are
\emph{damage}, \emph{attack}, \emph{opponent's}, \emph{energy} and
\emph{flip} (Figure~\ref{fig:terms}); the top bigrams (\emph{flip coin},
\emph{energy attached}, \emph{search deck}, \emph{damage counters}) are
game mechanics. Keyword patterns show how often key mechanics appear:
coin flips 26.9\%, Special Conditions 18.2\%, deck search 8.8\%, discard
pile 5.9\%, drawing 5.2\%, Energy attachment/movement 5.2\%. Each mechanic
targeted by our information needs is therefore found on several hundred
cards, enough for a non-trivial ranking problem. The illustrator
distribution is long-tailed: 5ban Graphics (946 cards), Ken Sugimori (834),
Kouki Saitou (601) and Mitsuhiro Arita (542) lead among 427 artists.
\todo{optional: NER over trivia/release info (e.g.\ spaCy) for events and
products}

\begin{figure}[t]
  \centering
  \figorbox{figures/damage_over_time.pdf}{\columnwidth}
  \caption{Median and 90th-percentile attack damage per release year.}
  \label{fig:dmg}
\end{figure}

\begin{figure}[t]
  \centering
  \figorbox{figures/top_terms.pdf}{\columnwidth}
  \caption{Most frequent terms in card rules text.}
  \label{fig:terms}
\end{figure}

% ==================================================
% 5. PROSPECTIVE SEARCH TASKS
% ==================================================

\section{Prospective Search Tasks}

\subsection{Description}

The search system retrieves and ranks card documents that satisfy a user's
information need, rather than relying only on exact keyword matches. A
relevant document provides information that addresses the user's
underlying need; a document that contains query terms in an unrelated
context is not relevant. For example, a card is relevant to a
discard-based strategy when its attacks or abilities interact with the
discard pile, not when the phrase appears incidentally. Relevance will be
judged on the card's rules text together with the enriched text from linked
entities.

\subsection{Information Needs}

\paragraph{IN1: Cards with discard-pile interactions.}
A player exploring discard-based strategies wants cards whose attacks or
abilities use or recover cards from the discard pile. Query:
\textit{``attacks discard pile''}. Relevant cards have effects that depend
on or retrieve discarded cards, not ones that merely discard as a cost.

\paragraph{IN2: Powerful attacks with specific Energy requirements.}
A player wants Pok\'emon with high-damage attacks that use a particular
Energy type. Query: \textit{``high damage fire energy attack''}. Relevance
combines the structured attack cost and damage with the effect text (a
``30+'' attack can be stronger than a fixed 60), so it cannot be answered
by a single filter.

\paragraph{IN3: Cards depicting Pok\'emon from a given habitat.}
A collector building a themed binder wants cards showing Pok\'emon that
live in the sea. Query: \textit{``Pok\'emon that live in the ocean''}.
Relevance comes from the linked species biology and Pok\'edex entries
rather than the card text, which tests the value of entity linking.

\paragraph{IN4: Cards that draw or search for cards.}
A player wants cards whose effects draw cards or search the deck. Query:
\textit{``ability draw cards search deck''}. Many cards are relevant, some
much more than others (e.g.\ repeatable abilities versus a one-off attack).

\paragraph{IN5: Cards by illustrator.}
A collector wants cards illustrated by a particular artist. Query:
\textit{``cards illustrated by Mitsuhiro Arita''}. Illustrator metadata lets
the system match cards whose own text never names the artist.

% ==================================================
% 6. CONCLUSIONS
% ==================================================

\section{Conclusions}

We built a collection of 13{,}972 Pok\'emon TCG cards from Bulbapedia,
linked to species, illustrators and expansions and enriched with TCGplayer
prices, and characterized its completeness, evolution over time and
vocabulary. Its main strengths are the breadth of its coverage (1999--2026)
and the amount of text linked to each card. Its main limitations are that basic Energy cards
are missing, 22\% of TCGCSV products could not be matched, and part of the
expansions lack a parsed release year. Next steps are the cleaning step
described in Section~3, building the denormalized documents and indexing
them in Solr for Milestone~2. Competitive data (decklists from tournament
APIs such as Limitless~\cite{limitless}) is a candidate extension in
Milestone~3.

% ==================================================
% REFERENCES
% ==================================================

\bibliographystyle{ACM-Reference-Format}
\bibliography{references}

\end{document}
